\section*{Supporting information}

\paragraph*{S1 Fig.} \textbf{Site frequency spectrum distributions per super-population and variant type.} Detailed SFS histograms showing the distribution of minor allele frequencies for SNVs, INDELs, and SVs within each of the five super-populations (AFR, EUR, EAS, SAS, AMR). SVs are dramatically skewed toward rare variants across all populations.

\paragraph*{S2 Fig.} \textbf{Count-matched PCA concordance.} PCA embeddings when all three variant types are subsampled to 18,519 variants (the SV count), demonstrating that SNV-INDEL concordance ($r \approx 0.65$) exceeds SNV-SV concordance ($r \approx 0.55$) even at matched variant counts.

\paragraph*{S3 Fig.} \textbf{SV subtype site frequency spectra.} Rare-variant proportion (MAF $< 0.05$) for deletions and insertions across five super-populations, showing that insertions are consistently more rare across all populations.

\paragraph*{S4 Fig.} \textbf{Variant size distributions.} Histograms of INDEL sizes (1--50 bp, linear scale) and SV sizes (50--200+ bp, log scale), confirming that the SV call set is dominated by variants in the 50--200 bp range.

\paragraph*{S5 Fig.} \textbf{Uniform MAF filter SFS comparison.} Site frequency spectra computed with a uniform MAF $\geq 0.01$ filter applied to all three variant types, confirming that the SV rare-variant enrichment (59--63\% with MAF 0.01--0.05) is biological rather than an artifact of differential MAF thresholds.
